% TODO proof-reading
\subsection{値を返す}

おそらく最もシンプルなJava関数は、何らかの値を返すものです。

ああ、そしてJavaには一般的な意味での「フリー」関数は存在しないことを心に留めておく必要があります。
それらは「メソッド」です。

各メソッドはあるクラスに関連しているため、クラスの外部でメソッドを定義することはできません。

しかし、簡単にするために、それらを「関数」と呼ぶことにします。

\begin{lstlisting}[style=customjava]
public class ret
{
	public static int main(String[] args) 
	{
		return 0;
	}
}
\end{lstlisting}

コンパイルしてみましょう：

\begin{lstlisting}
javac ret.java
\end{lstlisting}

\dots そして標準のJavaユーティリティを使用して逆コンパイルします：

\begin{lstlisting}
javap -c -verbose ret.class
\end{lstlisting}

次のような結果が得られます：

\begin{lstlisting}[caption=JDK 1.7 (抜粋)]
  public static int main(java.lang.String[]);
    flags: ACC_PUBLIC, ACC_STATIC
    Code:
      stack=1, locals=1, args_size=1
         0: iconst_0      
         1: ireturn       
\end{lstlisting}

Java開発者は、0がプログラミングで最も忙しい定数の一つであると判断し、
0をプッシュする専用の短い1バイトの\GTT{iconst\_0}命令があります

\footnote{MIPSでゼロ定数用の専用レジスタが存在するのと同様です：\myref{MIPS_zero_register}。}。
\GTT{iconst\_1}（1をプッシュ）、\GTT{iconst\_2}などもあり、\GTT{iconst\_5}まであります。

-1をプッシュする\GTT{iconst\_m1}もあります。

JVMでは、スタックが呼び出される関数にデータを渡すため、また戻り値のためにも使用されます。
したがって、\GTT{iconst\_0}は0をスタックにプッシュします。
\GTT{ireturn}は\ac{TOS}から整数値（\IT{i}は\IT{integer}を意味する）を返します。

例を少し書き直して、今度は1234を返してみましょう：

\begin{lstlisting}[style=customjava]
public class ret
{
	public static int main(String[] args)
	{
		return 1234;
	}
}
\end{lstlisting}

\dots 次のような結果が得られます：

\begin{lstlisting}[caption=JDK 1.7 (抜粋)]
  public static int main(java.lang.String[]);
    flags: ACC_PUBLIC, ACC_STATIC
    Code:
      stack=1, locals=1, args_size=1
         0: sipush        1234
         3: ireturn       
\end{lstlisting}

\GTT{sipush}（\IT{short integer}）は1234をスタックにプッシュします。
名前の\IT{short}は16ビット値がプッシュされることを示します。
数字1234は実際に16ビット値によく収まります。

より大きな値はどうでしょうか？

\begin{lstlisting}[style=customjava]
public class ret
{
	public static int main(String[] args) 
	{
		return 12345678;
	}
}
\end{lstlisting}

\begin{lstlisting}[caption=定数プール]
...
   #2 = Integer            12345678
...
\end{lstlisting}

\begin{lstlisting}
  public static int main(java.lang.String[]);
    flags: ACC_PUBLIC, ACC_STATIC
    Code:
      stack=1, locals=1, args_size=1
         0: ldc           #2                  // int 12345678
         2: ireturn       
\end{lstlisting}

JVM命令のオペコードに32ビット数をエンコードすることはできません。
開発者はそのような可能性を残しませんでした。

したがって、32ビット数12345678はいわゆる「定数プール」に格納されます。これは、
最もよく使用される定数（文字列、オブジェクトなどを含む）のライブラリのようなものです。

この定数を渡す方法はJVM独特のものではありません。

MIPS、ARMおよび他のRISC CPUも32ビット命令に32ビット数をエンコードできないため、
RISCのCPUコード（MIPSとARMを含む）は値を複数のステップで構築するか、
データセグメントに保持する必要があります：
\myref{ARM_big_constants}、\myref{MIPS_big_constants}。

MIPSコードも伝統的に「リテラルプール」と呼ばれる定数プールを持ち、
セグメントは「.lit4」（32ビット単精度浮動小数点定数用）と「.lit8」
（64ビット倍精度浮動小数点定数用）と呼ばれます。

他のデータ型も試してみましょう！

ブール値：

\begin{lstlisting}[style=customjava]
public class ret
{
	public static boolean main(String[] args) 
	{
		return true;
	}
}
\end{lstlisting}

\begin{lstlisting}
  public static boolean main(java.lang.String[]);
    flags: ACC_PUBLIC, ACC_STATIC
    Code:
      stack=1, locals=1, args_size=1
         0: iconst_1      
         1: ireturn       
\end{lstlisting}

このJVMバイトコードは、整数1を返すものと変わりません。

\CCppのように、ブール値にもスタック内の32ビットデータスロットが使用されます。

しかし、返されたブール値を整数として使用したり、その逆はできません --- 型情報はクラスファイルに格納され、実行時にチェックされます。

16ビット\IT{short}でも同じです：

\begin{lstlisting}[style=customjava]
public class ret
{
	public static short main(String[] args) 
	{
		return 1234;
	}
}
\end{lstlisting}

\begin{lstlisting}
  public static short main(java.lang.String[]);
    flags: ACC_PUBLIC, ACC_STATIC
    Code:
      stack=1, locals=1, args_size=1
         0: sipush        1234
         3: ireturn       
\end{lstlisting}

\dots そして\IT{char}も！

\begin{lstlisting}[style=customjava]
public class ret
{
	public static char main(String[] args) 
	{
		return 'A';
	}
}
\end{lstlisting}

\begin{lstlisting}
  public static char main(java.lang.String[]);
    flags: ACC_PUBLIC, ACC_STATIC
    Code:
      stack=1, locals=1, args_size=1
         0: bipush        65
         2: ireturn       
\end{lstlisting}

\INS{bipush}は「バイトをプッシュ」を意味します。
言うまでもなく、Javaの\IT{char}は16ビットUTF-16文字で、
\IT{short}と同等ですが、「A」文字のASCIIコードは65で、
スタックにバイトをプッシュする命令を使用することが可能です。

\IT{byte}も試してみましょう：

\begin{lstlisting}[style=customjava]
public class retc
{
	public static byte main(String[] args) 
	{
		return 123;
	}
}
\end{lstlisting}

\begin{lstlisting}
  public static byte main(java.lang.String[]);
    flags: ACC_PUBLIC, ACC_STATIC
    Code:
      stack=1, locals=1, args_size=1
         0: bipush        123
         2: ireturn       
\end{lstlisting}

なぜ内部的に32ビット整数として動作する16ビット\IT{short}データ型を使うのでしょうか？

\IT{short}データ型と同じなのに、なぜ\IT{char}データ型を使うのでしょうか？

答えは簡単です：データ型の制御とソースコードの可読性のためです。

\IT{char}は本質的に\IT{short}と同じかもしれませんが、それがUTF-16文字のプレースホルダーであり、
他の整数値のためのものではないことを素早く理解できます。

\IT{short}を使用する場合、変数の範囲が16ビットに制限されていることを皆に示します。

C スタイルの\IT{int}の代わりに、必要な場所で\IT{boolean}型を使用するのは非常に良いアイデアです。

Javaには64ビット整数データ型もあります：

\begin{lstlisting}[style=customjava]
public class ret3
{
	public static long main(String[] args)
	{
		return 1234567890123456789L;
	}
}
\end{lstlisting}

\begin{lstlisting}[caption=定数プール]
...
   #2 = Long               1234567890123456789l
...
\end{lstlisting}

\begin{lstlisting}
  public static long main(java.lang.String[]);
    flags: ACC_PUBLIC, ACC_STATIC
    Code:
      stack=2, locals=1, args_size=1
         0: ldc2_w        #2                  // long 1234567890123456789l
         3: lreturn       
\end{lstlisting}

64ビット数も定数プールに格納され、\INS{ldc2\_w}がそれをロードし、\INS{lreturn}
（\IT{long return}）がそれを返します。

\INS{ldc2\_w}命令は、定数プールから倍精度浮動小数点数
（これも64ビットを占める）をロードするためにも使用されます：

\begin{lstlisting}[style=customjava]
public class ret
{
	public static double main(String[] args)
	{
		return 123.456d;
	}
}
\end{lstlisting}

\begin{lstlisting}[caption=定数プール]
...
   #2 = Double             123.456d
...
\end{lstlisting}

\begin{lstlisting}
  public static double main(java.lang.String[]);
    flags: ACC_PUBLIC, ACC_STATIC
    Code:
      stack=2, locals=1, args_size=1
         0: ldc2_w        #2                  // double 123.456d
         3: dreturn       
\end{lstlisting}

\INS{dreturn}は「doubleを返す」を意味します。

そして最後に、単精度浮動小数点数：

\begin{lstlisting}[style=customjava]
public class ret
{
	public static float main(String[] args)
	{
		return 123.456f;
	}
}
\end{lstlisting}

\begin{lstlisting}[caption=定数プール]
...
   #2 = Float              123.456f
...
\end{lstlisting}

\begin{lstlisting}
  public static float main(java.lang.String[]);
    flags: ACC_PUBLIC, ACC_STATIC
    Code:
      stack=1, locals=1, args_size=1
         0: ldc           #2                  // float 123.456f
         2: freturn       
\end{lstlisting}

ここで使用される\INS{ldc}命令は、定数プールから32ビット整数をロードするのと同じものです。

\INS{freturn}は「floatを返す」を意味します。

それでは、何も返さない関数はどうでしょうか？

\begin{lstlisting}[style=customjava]
public class ret
{
	public static void main(String[] args) 
	{
		return;
	}
}
\end{lstlisting}

\begin{lstlisting}
  public static void main(java.lang.String[]);
    flags: ACC_PUBLIC, ACC_STATIC
    Code:
      stack=0, locals=1, args_size=1
         0: return        
\end{lstlisting}

これは、\INS{return}命令が実際の値を返すことなく制御を返すために使用されることを意味します。

これらすべてを知っていれば、最後の命令から関数（またはメソッド）の戻り型を
非常に簡単に推測することができます。